
% ════════════════════════════════════════════
% 具体分析（5.3.1.具体分析.tex，三级子块）—— 写作提示词见本文件末尾【本章节写作提示词】。
% ════════════════════════════════════════════
\subsubsection{具体分析}

针对问题三，本问属于误差传播与概率约束反解类问题，采用线性误差传播解析建模与蒙特卡洛仿真验证相结合的方法进行求解。

首先由车辆定位误差与方位角的几何关系建立指令方位角偏差的线性化模型，接着将三类独立误差源的标准差合成得到总方位偏差标准差，然后按正态分布 95\% 分位数反解定位标准差上限与最大允许行驶距离，最后通过 20 万次蒙特卡洛验证小角度近似的适用性并比较三类误差源贡献。

具体分析流程如图\ref{fig:analysis3_flow}所示。

\begin{figure}[ht]
    \begin{center}
        \begin{tikzpicture}[x=1cm, y=0.7cm]
            \node[box] (n11) at (0, 0) {定位误差建模};
            \node[box] (n12) at (5, 0) {方位偏差线性化};
            \node[box] (n13) at (10, 0) {三误差源合成};
            \node[box] (n22) at (5, -3) {正态分位数};
            \node[box] (n21) at (0, -3) {概率约束反解};
            \node[box] (n23) at (10, -3) {参数约束推导};
            \node[box] (n31) at (0, -6) {蒙特卡洛验证};
            \node[box] (n32) at (5, -6) {方差分解};
            \node[box] (n33) at (10, -6) {部署建议};
            \draw[arrow] (n11) -- (n12);
            \draw[arrow] (n12) -- (n13);
            \draw[arrow] (n23) -- (n22);
            \draw[arrow] (n22) -- (n21);
            \draw[arrow] (n13) -- ++(0,-1.0) -| (n23);
            \draw[arrow] (n21) -- (n31);
            \draw[arrow] (n31) -- (n32);
            \draw[arrow] (n32) -- (n33);
        \end{tikzpicture}
        \caption{问题三分析流程图}
        \label{fig:analysis3_flow}
    \end{center}
\end{figure}
